\section{Sistemos kūrimo projektas}

Šiame skyriuje pateikiama sistemos reikalavimų specifikacija, projektavimo metodai, statinis ir dinaminis sistemos vaizdas.

\subsection{Sistemos reikalavimų specifikacija}

\subsubsection{Komercinė specifikacija}

\lipsum[1]

\subsubsection{Sistemos funkcijos}

Funkciniai reikalavimai išreiškiami panaudojimo atvejų (toliau -- PA) modeliu. Kiekvienas PA aprašomas specifikacijos lentele, kurią sukuria makrokomanda \verb|\usecasespec| (žr. \ref{tab:pa-perziureti-irasus} lentelę). Reikšmes, kuriose yra kablelių, būtina apgaubti riestiniais skliaustais.

\usecasespec{
  label=tab:pa-perziureti-irasus,
  pa=Peržiūrėti įrašus,
  tikslas=Atvaizduoti visus naudotojui prieinamus įrašus.,
  aprasymas={Prisijungęs naudotojas pasirenka meniu punktą „Įrašai“. Sistema pateikia įrašų sąrašą, kurį galima filtruoti ir rikiuoti.},
  pries=Naudotojas prisijungęs prie sistemos.,
  aktoriai={Naudotojas, Administratorius},
  ispleciantys=Peržiūrėti įrašo informaciją,
  apimami=--,
  specializuojami=--,
  po=Naudotojas mato įrašų sąrašą.
}

\subsubsection{Naudotojo sąsajos specifikacija}

\lipsum[2]

\subsubsection{Realizacijai keliami reikalavimai}

\lipsum[3]

\subsubsection{Techninė specifikacija}

\lipsum[4]

\subsection{Sistemos projektavimo metodai}

\subsubsection{Projektavimo valdymas ir eiga}

\lipsum[5]

\subsubsection{Projektavimo technologija}

\lipsum[6]

\subsubsection{Programavimo kalbos, derinimo, automatizavimo priemonės, operacinės sistemos}

\lipsum[7]

\subsection{Sistemos projektas}

\subsubsection{Statinis sistemos vaizdas}

\paragraph{Diegimo diagrama.} \lipsum[8][1-4]

\paragraph{Komponentų diagrama.} \lipsum[9][1-4]

\subsubsection{Dinaminis sistemos vaizdas}

\lipsum[10]
